% Generated draft: requires mathematical and source review.
\documentclass[11pt]{article}
\usepackage[margin=25mm]{geometry}
\usepackage{fontspec}
\setmainfont{Latin Modern Roman}
\usepackage{amsmath,amssymb,mathtools,unicode-math}
\setmathfont{Latin Modern Math}
\usepackage{enumitem,needspace,xurl,hyperref}
\setlength{\parindent}{0pt}
\setlength{\parskip}{5pt}
\setlength{\emergencystretch}{3em}
\newcommand{\N}{\mathbb{N}}
\newcommand{\Z}{\mathbb{Z}}
\newcommand{\Q}{\mathbb{Q}}
\newcommand{\R}{\mathbb{R}}
\newcommand{\C}{\mathbb{C}}
\newcommand{\F}{\mathbb{F}}
\newcommand{\A}{\mathbb{A}}
\newcommand{\PP}{\mathbb{P}}
\newcommand{\E}{\mathbb{E}}
\newcommand{\eps}{\varepsilon}
\newcommand{\dd}{\,\mathrm{d}}
\begin{document}
\section*{UnsolvedMath ID: 700281}
\textbf{Source problem code:} catalogue entry
\subsection{Definitions and mathematical statement}
A \emph{\(2\)-group} is a group \(G\) such that every element has finite order equal to a power of \(2\); the group itself may be infinite. A group is \emph{simple} if it is nontrivial and its only normal subgroups are \(\{1\}\) and the whole group.

For the formulation below, interpret \emph{binary-finite} as meaning that every subgroup generated by two elements is finite:
\[
\forall x,y\in G,\qquad |\langle x,y\rangle|<\infty.
\]
The supplied record does not define this term, so this interpretation requires confirmation against the source. It does not impose finiteness on every finitely generated subgroup.

\textbf{Part (b), the remaining question in the supplied record.} Is it true that every binary-finite \(2\)-group \(G\) with \(|G|>2\) has a nontrivial proper normal subgroup? Equivalently, under the interpretation above, is the following implication true for every group \(G\)?
\[
\begin{gathered}
|G|>2,\qquad
\forall g\in G\;\exists k\in\Z_{\geq 0}\quad \operatorname{ord}(g)=2^k,\\
\forall x,y\in G\quad |\langle x,y\rangle|<\infty
\end{gathered}
\quad\Longrightarrow\quad
\exists N\trianglelefteq G\quad \{1\}\subsetneq N\subsetneq G.
\]

\textbf{Part (a), supplied as archived and solved.} The original question asks whether every \(2\)-group of order greater than \(2\) is non-simple, with no binary-finiteness hypothesis. The supplied archive reports a negative answer, asserting the existence of infinite simple \(2\)-groups. It attributes this existence to results on finitely generated infinite groups of exponent \(2^n\), together with the positive solution of the Restricted Burnside Problem for exponent \(2^n\). Here exponent \(2^n\) means that \(g^{2^n}=1\) for every group element. This is a report of the supplied archive assessment, not an independently verified result or a proof in this draft.
\subsection{Short English statement}
Part (b) asks whether every \(2\)-group with more than two elements in which every two elements generate a finite subgroup must have a nontrivial proper normal subgroup. The meaning of ``binary-finite'' requires source confirmation. The supplied record labels the problem partially solved and reports a negative answer to part (a), which omits the binary-finiteness hypothesis.
\subsection{Sources}
\begin{itemize}
\item Local UnsolvedMath registry: supplied complete record with permanent ID \(700281\), exact source code \texttt{catalogue entry}, titled \emph{Kourovka Notebook Problem 4.74}. Proposed by V. P. Shunkov in 1973.
\item \emph{The Kourovka Notebook (Unsolved Problems in Group Theory)}, 4th issue (1973), Problem 4.74; supplied record describes an update in the 21st issue (2026), arXiv:1401.0300v48. The collection description names E. I. Khukhro and V. D. Mazurov as editors. Part (b) is supplied as remaining in the main section; part (a) is supplied from the Archive of Solved Problems.
\item S. V. Ivanov, \emph{Int. J. Algebra Comput.} \textbf{4}, nos. 1--2 (1994), 1--308, cited in the supplied archive explanation concerning finitely generated infinite groups of exponent \(2^n\).
\item I. G. Lys\"enok, \emph{Izv. Math.} \textbf{60}, no. 3 (1996), 453--654, cited in the same supplied archive explanation.
\item E. I. Zel'manov, \emph{Math. USSR--Sb.} \textbf{72} (1992), 543--565, cited in the supplied archive explanation for the positive solution of the Restricted Burnside Problem for groups of exponent \(2^n\).
\item The supplied status review is dated \texttt{2026-10-16} and assigns the label \texttt{partially\_solved}, describing archived solved parts and remaining parts in the main section. This assessment is supplied dated curation, not independently verified. The status-review quotation is truncated; the fuller bibliography above comes from the supplied background field.
\end{itemize}
\end{document}
